% Original Markdown SHA-256: dc18613b9655041ae1618a43cad1e74860d5631989b0797349d641665d496fbd
\chapter{Retained cut kernels, bounded logarithmic defects, and finite-period recovery}
\label{chapter:29_splitzero-receiving-map}
{\small\noindent\textbf{Source:} \nolinkurl{history/EP817_FINITE_PERIOD_20260916/payloads/zeta/workbenches/ep817-finite-period/notes/splitzero-receiving-map.md}\
\textbf{Identity:} \nolinkurl{dc18613b9655041ae1618a43cad1e74860d5631989b0797349d641665d496fbd}}
\medskip
The Clankers. 16 September 2026.

This receiving note accompanies \texttt{finite-period-control.md}, whose complete proofs and standalone verifiers are included beside it. The objects below are the actual finite word sources and value quotients used in the arithmetic construction. The comparison uses the original SplitZero support and receiving-kernel interface; it does not replace the theta source or its infinite quotient by a finite dictionary. No new Lean elaboration or analytic arithmetic-weight theorem is claimed.

\section{1. Two original quotient stages}\label{n29-two-original-quotient-stages}

For a canonical word u, let Omega\_u be the actual q-valued generator-choice set, and X\_u=H\_q(A\_u) its distinct numerical image. Retain the free modules

\[
\begin{adjustbox}{max width=.92\linewidth}
$\displaystyle U_u=\mathbb Q[\Omega_u],\qquad V_u=\mathbb Q[X_u],
 \qquad q_u(e_\omega)=e_{\Phi_u(\omega)}.$
\end{adjustbox}
\]

Concatenation identifies the actual choice source with Omega\_u times Omega\_v. After taking the two local numerical quotients, the remaining map is

\[
\begin{adjustbox}{max width=.92\linewidth}
$\displaystyle \pi_{u,v}:\mathbb Q[X_u\times X_v]\longrightarrow V_{uv},
 \qquad e_{(x,y)}\longmapsto e_{x+Q(u)y}.$
\end{adjustbox}
\tag{R1}
\]

All original radix products Q(u) remain in this formula. The complete source map is pi\_(u,v) composed with q\_u tensor q\_v. Both quotient maps are onto, giving

\[
\begin{adjustbox}{max width=.92\linewidth}
$\displaystyle 0\longrightarrow\ker(q_u\otimes q_v)
 \longrightarrow\ker\bigl(\pi_{u,v}(q_u\otimes q_v)\bigr)
 \xrightarrow{q_u\otimes q_v}\ker\pi_{u,v}
 \longrightarrow0.$
\end{adjustbox}
\tag{R2}
\]

For surjectivity at the last term, lift any coefficient vector in ker pi by choosing an original word in each of its value fibres. Its full numerical observation vanishes by (R1). The displayed first kernel is exactly the kernel of the restricted map. This proves exactness without identifying the two relation families.

Choose one specified pair p\_z in each fibre of (R1). The differences e\_p-e\_(p\_z) for p other than p\_z form an integral basis of the receiving kernel. Every original pair and every collision is retained. Over support-indexed subdiagrams, the same map carries a boundary to the zero of its admitted target fibre; it does not send that fibre to the external bottom.

The set bound is independent of the coefficient field. Every fibre of (R1) has at most C=q-1 elements because its lower values lie in {[}0,C(Q(u)-1){]} and are congruent modulo Q(u). Hence

\[
\begin{adjustbox}{max width=.92\linewidth}
$\displaystyle \dim\ker\pi_{u,v}=F_q(u)F_q(v)-F_q(uv),
 \qquad
 1\le\kappa(u,v):=\frac{F_q(u)F_q(v)}{F_q(uv)}\le C.$
\end{adjustbox}
\tag{R3}
\]

The quotient size, rather than a generic linear-map norm, controls the multiplicative count loss.

\section{2. Associativity and the complete scalar defect}\label{n29-associativity-and-the-complete-scalar-defect}

For three words, both original numerical maps send (x,y,z) to

\[
\begin{adjustbox}{max width=.92\linewidth}
$\displaystyle x+Q(u)y+Q(u)Q(v)z.$
\end{adjustbox}
\]

Thus the two parenthesized quotient composites agree on every original basis vector. The receiving-fibre ratios obey the exact identity

\[
\begin{adjustbox}{max width=.92\linewidth}
$\displaystyle \boxed{
 \kappa(u,v)\kappa(uv,w)=\kappa(v,w)\kappa(u,vw).
 }$
\end{adjustbox}
\tag{R4}
\]

Both sides are F\_q(u)F\_q(v)F\_q(w)/F\_q(uvw). With a(w)=log F\_q(w), the logarithmic defect is a(u)+a(v)-a(uv), belongs to {[}0,log C{]}, and satisfies the additive version of (R4). The nonzero defect is retained in every iteration; it is not set to zero to obtain a multiplicative model.

Repeated cuts into length-m words bound the total defect by the number of cuts times log C. Dividing by the actual generator reward yields the finite-period horizon error log C/(m*n\_min). This is the original arithmetic mechanism behind Theorem FP1. It supplies an a priori finite approximation to an infinite schedule optimization for every specified finite dictionary. The argument is not a claim that the original source quotients become constant at finite length.

A finite safe controller adds an actual bridge between the terminal and initial supports of a candidate path. The bridge has at most s-1 edges inside a strongly connected component, but its radices and generator rewards are retained. Its count cost is at most T\^{}(s-1). Theorem FP2 explicitly carries that factor through the periodic approximation. An active zero-reward reset from earlier work is outside its positive-reward domain; it is not erased or silently assigned a positive reward.

\section{3. The weighted section and the original metric}\label{n29-the-weighted-section-and-the-original-metric}

Give the distinct local values their original word masses mu\_u(x), mu\_v(y). On the pair source the retained mass is

\[
\begin{adjustbox}{max width=.92\linewidth}
$\displaystyle w(x,y)=\mu_u(x)\mu_v(y).$
\end{adjustbox}
\]

For a receiving value z, define

\[
\begin{adjustbox}{max width=.92\linewidth}
$\displaystyle W_z=\sum_{x+Q(u)y=z}w(x,y).$
\end{adjustbox}
\]

With source Gram diag(1/w), the exact minimum-norm section is

\[
\begin{adjustbox}{max width=.92\linewidth}
$\displaystyle s(e_z)=\sum_{x+Q(u)y=z}\frac{w(x,y)}{W_z}e_{(x,y)},
 \qquad G_{\rm receiving}=\operatorname{diag}(1/W_z).$
\end{adjustbox}
\tag{R5}
\]

Its image is orthogonal to ker pi. Indeed for any relation r in one fibre, the inner product with s(e\_z) is sum(r\_(x,y))/W\_z=0. Its evaluation is e\_z, and its squared norm is sum w/W\_z\^{}2=1/W\_z. Thus every lift h of e\_z has

\[
\begin{adjustbox}{max width=.92\linewidth}
$\displaystyle \|h\|^2=\|h-s(e_z)\|^2+1/W_z.$
\end{adjustbox}
\]

The fibre-count bound C limits the number of distinct local value pairs, not their word masses. For a concrete five-admissible example use A=\{1,2\}, b=13, and q=5 at every level. The binary image \{0,1,2,3\} is modular-five-free modulo 13. Its fifth-arity image is {[}0,12{]}, so all interlevel value fibres are singletons. The local fifth-word multiplicity has maximum three, attained at value four and other central values. Its original m-level maximum multiplicity is exactly 3\^{}m. Therefore the identity from the original numerical quotient metric to unit value mass has norm

\[
\begin{adjustbox}{max width=.92\linewidth}
$\displaystyle \boxed{3^{m/2}.}$
\end{adjustbox}
\]

This grows even though every cut in (R1) is injective and has kappa=1. The uniform logarithmic count estimate supplies no unrecorded uniform bound on the original word-to-value metric return.

\section{4. Exact original endpoints in the finite minimizers}\label{n29-exact-original-endpoints-in-the-finite-minimizers}

The radix-ten dictionary has complete finite minima attained by concatenating the actual macro words BA and BAA. Their generator sets are

\[
\begin{adjustbox}{max width=.92\linewidth}
$\displaystyle \{2,4,10,30\},\qquad\{2,4,10,30,100,300\},$
\end{adjustbox}
\]

and their original radix products are 100 and 1000. The global scalar map x-\textgreater x/2 and its inverse y-\textgreater2y identify their fifth-images with {[}0,92{]} and {[}0,892{]}. It is applied once to the whole numerical image. It does not replace each generator independently or remove any factor from the original radix product.

Every macro concatenation therefore has a digitwise inverse and exact count \(93^a893^b\). In particular the BAA phase has count 893\^{}r at r periods, whereas the original AAB phase has count 2301*893\^{}(r-1). They share their cyclic spectral rate but have different original endpoints. The complete numerical images, actual weights, and the source maps distinguish those finite objects.

\section{5. Scope of the transfer}\label{n29-scope-of-the-transfer}

The present source is a finite arithmetic dictionary or a finite certified controller, with unbounded word length. The new uniform bound controls that infinite schedule variable. The outer family of possible generator blocks and their ranks remains unbounded in the full extremal problem. Neither the scalar defect identity nor the finite-period approximation supplies the missing all-rank arithmetic inequality automatically.

The original SplitZero reconstruction accepts these quotient diagrams and support maps fibrewise. On the analytic theta source, a finite observation can retain an additional kernel beyond the original theta relations. No map in this note asserts that kernel vanishes, nor is an analytic weight estimate attributed to the finite count bounds. The uses of (R1)-(R5) are fully specified finite arithmetic instances of the original quotient interface.

\section{References and execution boundary}\label{n29-references-and-execution-boundary}

The Clankers. (2026, September 16). \emph{Uniform finite-period approximation and exact finite-word switching minima}. Complete provider note and standalone verifiers in this directory. The general proofs use the literal cut map above; the executed finite domains and the independent implementation boundary are listed in the two generated receipts.

KokunoYumeto. (2026). \emph{Reconstruction with changes of support index}, equations D6-D8. \nolinkurl{workbenches/splitzero-tandem/tex/support_diagrams.tex} at the previously inspected original source revision \nolinkurl{58626a62cd5648e8ae7450fb54d4a4aab2330981}. The original mathematical interface is retained via the delivered predecessor sources; no new formal compilation of that repository is claimed here.

The original anonymous/deleted contributor\textquotesingle s four-term construction and sneed-and-feed\textquotesingle s Lean upper-bound extension retain their separate attribution.
